\scanpage{14}
\uline{Démonstration} \quad $A,A'$ et l'idéal $\unsure{m|\mathfrak m}$ satisfont les conditions de prop 1, de plus $A_0^{\unsure{2|\deleted{2}}}=A/m=k$ est un corps, et $A'/mA'=A'/m\simeq A'_0$ est un vectoriel de dim 2 sur $k$, il est en particulier fid plat sur $k$. Alors (ii) et (iii) résultent du N\textsuperscript{o} 1, Prop. 1, cor. 5, et prop. 2 \inserted{cor. 1 et} cor. 3.

Pour (i), utilisant N\textsuperscript{o} 1, Prop 3, cor. 2, on est ramené au cas où $A=k$, donc $A'$ est une algèbre de rang 2 sur $k$. \inserted{On peut supposer $k$ alg. clos.} On doit distinguer deux cas

(i) \inserted{$A'$ réduit, d'où :} $A'\simeq k\times k$, alors la question est purement ensembliste ; il y a exactement 2 relations d'équivalence sur un ens à 2 éléments.

(ii) $A'\deleted{\simeq k}$ non réduit, alors $A'\simeq k[\varepsilon]$, avec $\varepsilon^2=0$. (\unsure{nbs|nbrs|nbres} duaux).
% archive sheet number: 13
